%
% writer_publish_generation.tex
%
% A Z specification of the Hub-writer publish path's inbox delivery --
% src/punt_lux/domain/hub/inbox.py's `_writer` closure (registered by
% `ensure_writer`, invoked by `Hub.publish`'s fan-out) -- prompted by a
% Copilot review finding on PR #498 (bead lux-m3xr) against
% inbox.py:92-93.
%
% WHY THIS MODEL EXISTS. `docs/menu_lifecycle.tex` proved D1: `offer`'s
% get-then-put across a released `_inboxes_lock` hold lets a `drop_session`
% land in the window and produce a false "delivered" report. That fix --
% get and put under one lock hold -- was applied to `offer` only. The
% session writer's own put,
%
%   def _writer(message: ObserverMessage) -> None:
%       inbox_for(connection_id).put(message)
%
% (inbox.py:92-93) still calls `inbox_for` -- which takes and releases
% `_inboxes_lock` internally -- and THEN calls `.put` on the object it
% returned, outside any lock. `offer`'s own docstring already names this
% as a different case ("Unlike the session writer's
% `inbox_for(...).put` -- which creates an inbox on demand -- this reads
% with `.get`..."); it does not claim the writer's put shares `offer`'s
% atomicity fix, and it does not.
%
% This model asks a sharper question than D1's presence-only check: not
% just whether the target is PRESENT when the put lands, but whether it
% is the connection's CURRENT inbox. `connection_lease_reaping.tex`'s I6
% and I7 establish that a live writer keeps its registry-level BINDING
% (`hasWriter`) and that departure clears all three cascade-leg flags,
% but that model's own commentary is explicit that it "has no way to
% distinguish a fresh writer from a stale one [by identity]" -- it
% tracks a boolean per connection, never an inbox's generation. Neither
% `menu_lifecycle.tex`'s M1 (offer only) nor `menu_lifecycle_delivery.tex`'s
% M1b (offer only, live-recipient loss) reaches the writer's put at all.
% This model closes that gap: it gives each connection's inbox a
% GENERATION -- the object identity a fresh `BoundedInbox()` gets on
% every (re)creation -- and proves a writer's put always lands in the
% generation live at PUT time, not the one live when it started looking.
%
% Read first: src/punt_lux/domain/hub/inbox.py (`inbox_for`, `_writer`,
% `ensure_writer`, `drop_session`, `offer`), docs/menu_lifecycle.tex (D1,
% the offer-path analogue this refines by scope), docs/connection_lease_reaping.tex
% (I6/I7, the coarse `hasWriter`/`inboxNonEmpty` boolean this sharpens
% into an object generation).
%
% Type-check:  fuzz docs/writer_publish_generation.tex
% Model-check: probcli docs/writer_publish_generation.tex -model_check -p DEFAULT_SETSIZE 3
% WG is a reachability check; see the Verification section for the exact
% -goal predicates.
%

\documentclass[a4paper,10pt,fleqn]{article}
\usepackage[margin=1in]{geometry}
\usepackage{fuzz}
\usepackage[colorlinks=true,linkcolor=blue,citecolor=blue,urlcolor=blue]{hyperref}

\begin{document}

\title{lux Hub-Writer Publish Path: a Z Specification}
\author{Formal model of the session writer's inbox put -- proving a
  publish always lands in the connection's currently-live inbox, never a
  detached predecessor or an unrelated successor}
\date{September 2026}
\maketitle

% ============================================================
\section{Introduction}
% ============================================================

Every MCP session with a live inbox has a writer closure registered
against its \texttt{ConnectionId} in the Hub's per-connection writer
registry. \texttt{Hub.publish} fans a message out by calling each
subscriber's writer; for the inbox writer, that call is
\texttt{inbox\_for(connection\_id).put(message)} -- look the inbox up
(creating it if none exists), then put into whatever was returned.

A connection's inbox is not fixed for the life of the \texttt{ConnectionId}:
\texttt{drop\_session} removes it on departure, and a later
\texttt{ensure\_writer} call -- on a same-identity reconnect, a scenario
\texttt{connection\_lease\_reaping.tex} already treats as a first-class
case -- creates a brand new one. Two calls a millisecond apart can
therefore be talking about two different inbox objects for the identical
\texttt{ConnectionId}. This model names that fact a \emph{generation}: the
inbox a connection has right now, distinct from any inbox it had before or
will have after.

\begin{description}
  \item[WG --- a publish lands only in the live generation.] A writer's put
    is committed against the inbox it is the CURRENT one for a connection at
    the moment the put actually happens, not the moment the writer began
    resolving which inbox to use. Violating this has two faces, both closed
    by the same property: the target has since been dropped and never
    recreated (the message is written into a detached, unreachable object --
    silent loss), or the target has since been dropped and recreated for a
    reconnected session (the message is written into a \emph{different, live}
    session's inbox -- delivery to the wrong recipient, not merely loss).
\end{description}

WG is deliberately kept out of the state schema, for the reason
\texttt{connection\_lease\_reaping.tex} and \texttt{menu\_lifecycle.tex}
both give: a property enforced as a guard on every successor would make the
violating operation unreachable rather than a counterexample. The schema
below carries only structural coherence; WG is checked as the reachability
of its negation (Section~\ref{sec:verify}).

Five operations realize the fixed design. \texttt{Arm} models
\texttt{inbox\_for} creating (or finding) a connection's inbox, called
standalone by a reconnecting \texttt{ensure\_writer}. \texttt{Depart} models
\texttt{drop\_session}. \texttt{WriterPut} models \texttt{\_writer}
\emph{once fixed}: the lookup-or-create and the put are one atomic step,
mirroring \texttt{offer}'s D1 fix of holding \texttt{\_inboxes\_lock} across
both instead of releasing it in between.

% ============================================================
\section{Basic Types}
% ============================================================

\begin{zed}
  [CONN]
\end{zed}

A connection identity -- the \texttt{ConnectionId} keying the inbox
registry. One \texttt{CONN} value suffices to exhibit a self-race between
one connection's own stalled writer and its own departure-then-reconnect;
the carrier is bounded by ProB's \texttt{DEFAULT\_SETSIZE}.

% ============================================================
\section{Free Types}
% ============================================================

\begin{zed}
  GEN ::= gen0 | gen1
\end{zed}

The generation of a connection's inbox: which incarnation of
\texttt{BoundedInbox} is the one \texttt{\_inboxes[connection\_id]} names
right now. Two values suffice: one recreation is enough to show a stalled
writer's captured generation diverging from the live one. Prefixed
constructors, matching the convention every companion model in this repo
uses for a small enumerated carrier.

\begin{zed}
  FSTATE ::= fclear | fset
\end{zed}

A two-value flag, matching \texttt{connection\_lease\_reaping.tex}'s and
\texttt{menu\_lifecycle.tex}'s \texttt{FSTATE}. Here it records whether any
put has ever landed against a generation other than the one live when the
put actually happened.

\begin{zed}
  PPHASE ::= pidle | pput
\end{zed}

Whether a writer's put is between resolving which inbox to use and
committing into it. In this, the fixed model, \texttt{WriterPut} is atomic,
so this never leaves \texttt{pidle}; the split fidelity control drives it to
\texttt{pput}, the window the real \texttt{inbox\_for(...).put(...)} split
exposes.

% ============================================================
\section{State}
% ============================================================

One flat schema, structural coherence only. WG is checked as the
reachability of its negation (Section~\ref{sec:verify}), not carried here.

\begin{schema}{WriterReg}
  hasInbox : \power CONN \\
  inboxGen : CONN \fun GEN \\
  putPhase : PPHASE \\
  putFor : \power CONN \\
  capturedGen : CONN \fun GEN \\
  genMismatch : FSTATE
\where
  \# putFor \leq 1
\end{schema}

\texttt{hasInbox} is the set of connections with a live, present inbox --
\texttt{connection\_id in \_inboxes}. \texttt{inboxGen} is total over
\texttt{CONN} for the same reason \texttt{connection\_lease\_reaping.tex}
keeps \texttt{owner} total over \texttt{SCENE}: a partial function keyed
like this would force every operation that touches a connection to decide,
separately, whether it is being seen for the first time; a generation value
persisting for a connection currently \emph{outside} \texttt{hasInbox} is
exactly what lets \texttt{Depart} preserve it and a stalled writer's earlier
capture be compared against it later. \texttt{putPhase}/\texttt{putFor} are
the split-window device \texttt{menu\_lifecycle.tex} uses for
\texttt{offerPhase}/\texttt{offerFor}: a phase plus the at-most-one
connection currently inside that window, inert (\texttt{pidle}, empty)
throughout this fixed model. \texttt{capturedGen} holds, for a writer
mid-put, the generation it resolved \texttt{inbox\_for} to; like
\texttt{inboxGen} it is kept total for the same reason.

% ============================================================
\section{Initialization}
% ============================================================

\begin{schema}{Init}
  WriterReg'
\where
  hasInbox' = \emptyset \\
  inboxGen' = (\lambda c : CONN @ gen0) \\
  putPhase' = pidle \\
  putFor' = \emptyset \\
  capturedGen' = (\lambda c : CONN @ gen0) \\
  genMismatch' = fclear
\end{schema}

No connection has an inbox, no put in flight, nothing yet mismatched.

% ============================================================
\section{Operations}
% ============================================================

\subsection{Arm --- resolve (creating if absent) a connection's inbox}

Models \texttt{inbox\_for(connection\_id)} called on its own, as
\texttt{ensure\_writer} calls it on every (re)connect. If the connection has
no inbox, one is created and its generation is set to a value distinct from
whatever the connection's generation last was -- a fresh
\texttt{BoundedInbox()} instance is never the departed one, whether this is
the connection's first inbox or its third. If an inbox is already present,
\texttt{inbox\_for} returns the existing one unchanged.

\begin{schema}{Arm}
  \Delta WriterReg \\
  c? : CONN
\where
  putPhase' = putPhase \\
  putFor' = putFor \\
  capturedGen' = capturedGen \\
  genMismatch' = genMismatch \\
  c? \notin hasInbox \implies
    (hasInbox' = hasInbox \cup \{ c? \} \land
     (\exists g : GEN | g \neq inboxGen(c?) @
       inboxGen' = inboxGen \oplus \{ c? \mapsto g \})) \\
  c? \in hasInbox \implies
    (hasInbox' = hasInbox \land inboxGen' = inboxGen)
\end{schema}

\subsection{Depart --- drop a connection's inbox}

Models \texttt{drop\_session}: the inbox is removed from \texttt{hasInbox}.
Its generation value is left exactly as it was -- \texttt{drop\_session}
does not overwrite anything, it pops a dict entry -- so a generation
persisting after departure is not a modeling convenience, it is the literal
behavior: the dropped \texttt{BoundedInbox} object still exists, still holds
whatever generation it was created with, and any reference to it a stalled
caller is holding remains valid to call \texttt{.put} on.

\begin{schema}{Depart}
  \Delta WriterReg \\
  c? : CONN
\where
  c? \in hasInbox \\
  hasInbox' = hasInbox \setminus \{ c? \} \\
  inboxGen' = inboxGen \\
  putPhase' = putPhase \\
  putFor' = putFor \\
  capturedGen' = capturedGen \\
  genMismatch' = genMismatch
\end{schema}

\subsection{WriterPut --- resolve and put, one atomic step (the fix)}
\label{sec:writerput}

Models \texttt{\_writer} \emph{once fixed}: the shipped code calls
\texttt{inbox\_for(connection\_id)} -- which takes and releases
\texttt{\_inboxes\_lock} on its own -- and only then calls \texttt{.put} on
whatever it returned, outside any lock. Here the whole operation -- resolve
or create, then put -- is one atomic step, matching \texttt{offer}'s D1 fix
of holding \texttt{\_inboxes\_lock} across both instead of releasing it in
between. Because the resolution and the put are the same step, the put
always lands in the inbox this step itself just resolved; there is no
window for a \texttt{Depart} or a reconnecting \texttt{Arm} to change which
inbox that is before the put runs. The put itself carries no observable
message identity in this model -- like \texttt{connection\_lease\_reaping.tex}'s
\texttt{Activate}, "modeled abstractly," since WG needs only that a put
landed against the live generation, not which message or which topic --
so \texttt{WriterPut}'s effect on \texttt{hasInbox}/\texttt{inboxGen} is
identical to \texttt{Arm}'s; the operative difference the fix supplies is
that resolution and put are one step, which the split control below breaks
apart.

\begin{schema}{WriterPut}
  \Delta WriterReg \\
  c? : CONN
\where
  putPhase = pidle \\
  putPhase' = putPhase \\
  putFor' = putFor \\
  capturedGen' = capturedGen \\
  genMismatch' = genMismatch \\
  c? \notin hasInbox \implies
    (hasInbox' = hasInbox \cup \{ c? \} \land
     (\exists g : GEN | g \neq inboxGen(c?) @
       inboxGen' = inboxGen \oplus \{ c? \mapsto g \})) \\
  c? \in hasInbox \implies
    (hasInbox' = hasInbox \land inboxGen' = inboxGen)
\end{schema}

% ============================================================
\section{Invariant}
\label{sec:inv}
% ============================================================

\paragraph{WG --- a publish lands only in the live generation.}
\[
  \lnot\, (genMismatch = fset).
\]
No operation in this, the fixed, model ever sets \texttt{genMismatch}:
\texttt{Arm}, \texttt{Depart}, and \texttt{WriterPut} all carry
\texttt{genMismatch' = genMismatch} unconditionally. WG therefore holds by
inspection of every operation's predicate, not merely by the absence of a
witness ProB happens not to have found -- and Section~\ref{sec:verify}
confirms the same conclusion empirically. \texttt{WriterPut}'s atomicity is
the reason no operation \emph{needs} to touch \texttt{genMismatch}: because
resolution and put are the same step, the generation a put commits against
is always, trivially, the generation live at that instant. The split
fidelity control (Section~\ref{sec:fidelity}) shows what breaks when that
one step becomes two.

% ============================================================
\section{Fidelity Control}
\label{sec:fidelity}
% ============================================================

\paragraph{Split --- \texttt{writer\_publish\_generation\_split\_buggy.tex}.}
\texttt{WriterPut} is replaced by \texttt{WriterPutBegin}/\texttt{WriterPutEnd},
splitting the shipped code's \texttt{inbox\_for(connection\_id).put(message)}
at the exact point \texttt{inbox\_for} returns and releases
\texttt{\_inboxes\_lock}. \texttt{WriterPutBegin(c?)} resolves (creating if
absent, exactly as \texttt{Arm} does) and captures the resolved generation
into \texttt{capturedGen}, entering \texttt{pput}. \texttt{WriterPutEnd}
commits the put \emph{unconditionally} against whatever \texttt{capturedGen}
holds; only if the connection is \emph{still} present with that \emph{same}
generation does the put land correctly, and \texttt{genMismatch} is set the
moment it does not.

Two traces witness WG's negation, one for each face named in the
Introduction.

\emph{Orphan loss.} \texttt{Arm}$(c_1)$;
\texttt{WriterPutBegin}$(c_1)$; \texttt{Depart}$(c_1)$;
\texttt{WriterPutEnd} -- the departure runs in the window, and at
\texttt{WriterPutEnd} the connection is no longer in \texttt{hasInbox} at
all: the put lands in a detached object nothing will ever read again.

\emph{Cross-generation delivery -- the worse case.}
\texttt{Arm}$(c_1)$; \texttt{WriterPutBegin}$(c_1)$; \texttt{Depart}$(c_1)$;
\texttt{Arm}$(c_1)$; \texttt{WriterPutEnd}. The second \texttt{Arm} models a
same-identity reconnect's own \texttt{ensure\_writer} call, landing between
the stalled writer's resolution and its put -- exactly the scenario
\texttt{connection\_lease\_reaping.tex} already treats as ordinary
(\texttt{Connect}, "no precondition privileging a fresh connection over a
reconnecting one"). It creates a fresh inbox at the OTHER generation value.
At \texttt{WriterPutEnd} the connection IS present -- \texttt{hasInbox}
includes $c_1$ -- but \texttt{capturedGen}$(c_1)$ is the generation the
\emph{first} \texttt{Arm} set, not the one the second one just did. The put
commits anyway: a message published under the departed session's writer
lands in the reconnected session's live inbox, not lost, delivered to the
wrong recipient.

Both traces are \emph{FOUND} against the split control and \emph{NOT found}
against this, the fixed, spec, because \texttt{WriterPut}'s single step
gives a \texttt{Depart} or a reconnecting \texttt{Arm} no window to land in
between a resolution and a put that are the same transition.

% ============================================================
\section{Verification}
\label{sec:verify}
% ============================================================

Type-checking is a \texttt{fuzz} pass:
\begin{verbatim}
  fuzz docs/writer_publish_generation.tex
\end{verbatim}

WG is a safety property, checked by asking ProB whether its negation is
\emph{reachable}. A goal reported \emph{not found} means the invariant
holds on every reachable state.

\begin{verbatim}
  # WG (goal NOT found in this, the fixed, spec)
  probcli docs/writer_publish_generation.tex -model_check -nodead \
    -goal "genMismatch = fset" \
    -p DEFAULT_SETSIZE 2
\end{verbatim}

A correct put is reachable, so the model is not vacuously safe by never
putting at all:

\begin{verbatim}
  # A connection has a live, current inbox (goal FOUND)
  probcli docs/writer_publish_generation.tex -model_check -nodead \
    -goal "#c.(c : CONN & c : hasInbox)" \
    -p DEFAULT_SETSIZE 2
\end{verbatim}

Deadlock-freedom and the structural state invariant are checked by a full
model-check over the reachable state space:

\begin{verbatim}
  probcli docs/writer_publish_generation.tex -model_check -p DEFAULT_SETSIZE 2
\end{verbatim}

The fidelity control makes the negation \emph{FOUND}:

\begin{verbatim}
  # Split: WG's negation FOUND against the split control, NOT found against fixed
  probcli docs/writer_publish_generation_split_buggy.tex -model_check -nodead \
    -goal "genMismatch = fset" \
    -p DEFAULT_SETSIZE 2
\end{verbatim}

Re-run \texttt{fuzz} and these goal checks whenever \texttt{inbox.py}'s
\texttt{inbox\_for}, \texttt{ensure\_writer}, \texttt{\_writer}, or
\texttt{drop\_session} change.

\end{document}
